\newpage
\Article{69}
Viel schwieriger sind jedoch solche Moduln $a+bi$ zu behandeln, für die $b=0$ nicht gilt (Zahlen der vierten Art aus Art.~34); mehrere Untersuchungen müssen daher vorausgehen. Die Norm $aa+bb$, die eine reelle Primzahl der Form $4n+1$ sein wird, bezeichnen wir mit $p$.

Es sei $S$ die Sammlung aller einfach kleinsten Reste für den Modul $a+bi=m$ unter Ausschluss der Null, sodass die Anzahl der in $S$ enthaltenen Zahlen $=p-1$ ist. Es bezeichne $x+yi$ eine beliebige Zahl dieses Systems, und wir setzen $ax+by=\xi,ay-bx=\eta$. Dann sind $\xi,\eta$ ganze Zahlen zwischen den Grenzen $0$ und $p$, \emph{ohne die Grenzen selbst}. Denn im vorliegenden Fall, in dem $a,b$ teilerfremd sind, zeigen die Formeln aus Art.~45, nämlich $\eta\equiv k\xi,\xi\equiv-k\eta\pmod p$, dass keine der Zahlen $\xi,\eta$ $=0$ sein kann, ohne dass die andere zugleich verschwindet, sodass $x=0,y=0$ wird; diese Kombination haben wir bereits ausgeschlossen. Das Kriterium dafür, dass die Zahl $x+yi$ in $S$ enthalten ist, besteht also darin, dass die vier Zahlen $\xi,\eta,p-\xi,p-\eta$ positiv sind.

Ferner bemerken wir, dass für keine solche Zahl $\xi=\eta$ gelten kann. Daraus würde nämlich $p(x+y)=a(\xi+\eta)+b(\xi-\eta)=2a\xi$ folgen, was widersinnig ist, da keiner der Faktoren $2,a,\xi$ durch $p$ teilbar ist. Ebenso zeigt die Gleichung $p(x-y+a+b)=2a\xi+(a+b)(p-\xi-\eta)$, dass $\xi+\eta=p$ nicht gelten kann. Da die Zahlen $\xi-\eta,p-\xi-\eta$ daher entweder positiv oder negativ sein müssen, gewinnen wir daraus eine Unterteilung des Systems $S$ in vier Sammlungen $C,C',C'',C'''$, denen wir die Zahlen folgendermaßen zuordnen:
\[
\begin{array}{c|ll}
\text{In die Sammlung}&\multicolumn{2}{c}{\text{Zahlen, für die}}\\\hline
C&\xi-\eta\text{ positiv},&p-\xi-\eta\text{ positiv}\\
C'&\xi-\eta\text{ positiv},&p-\xi-\eta\text{ negativ}\\
C''&\xi-\eta\text{ negativ},&p-\xi-\eta\text{ negativ}\\
C'''&\xi-\eta\text{ negativ},&p-\xi-\eta\text{ positiv}
\end{array}
\]

Das Kriterium für eine Zahl der Sammlung $C$ ist also eigentlich sechsfach: Die sechs Zahlen $\xi,\eta,p-\xi,p-\eta,\xi-\eta,p-\xi-\eta$ müssen positiv sein. Doch die Bedingungen $2,5$ und $6$ schließen die übrigen offenbar bereits ein. Entsprechendes gilt für die Sammlungen $C',C'',C'''$, sodass die vollständigen Kriterien dreifach sind, nämlich:
\SourceLeaf{151}%
\[
\begin{array}{c|lll}
\text{Für die Sammlung}&\multicolumn{3}{c}{\text{Folgende Zahlen müssen positiv sein}}\\\hline
C&\eta,&\xi-\eta,&p-\xi-\eta\\
C'&p-\xi,&\xi-\eta,&\xi+\eta-p\\
C''&p-\eta,&\eta-\xi,&\xi+\eta-p\\
C'''&\xi,&\eta-\xi,&p-\xi-\eta
\end{array}
\]

\Needspace{10\baselineskip}
Auch ohne unseren Hinweis wird jeder leicht erkennen, dass bei der geometrischen Darstellung komplexer Zahlen (siehe Art.~39) die Zahlen des Systems $S$ im Inneren des Quadrats liegen, dessen Seiten die Punkte verbinden, die die Zahlen $0,a+bi,(1+i)(a+bi),i(a+bi)$ darstellen, und dass die Unterteilung des Systems $S$ der Teilung des Quadrats durch seine Diagonalen entspricht. Hier haben wir jedoch rein arithmetische Überlegungen vorgezogen und überlassen der Kürze halber die Veranschaulichung durch eine geometrische Darstellung dem erfahrenen Leser.

\Needspace{22\baselineskip}
\Article{70}
Sind vier komplexe Zahlen $r=x+yi,r'=x'+y'i,r''=x''+y''i,r'''=x'''+y'''i$ so miteinander verbunden, dass $r'=m+ir,r''=m+ir'=(1+i)m-r,r'''=m+ir''=im-ir$ gilt, und nimmt man an, dass die erste, $r$, zur Sammlung $C$ gehört, so gehören die übrigen $r',r'',r'''$ der Reihe nach zu den Sammlungen $C',C'',C'''$. Setzt man nämlich $\xi=ax+by,\eta=ay-bx,\xi'=ax'+by',\eta'=ay'-bx',\xi''=ax''+by'',\eta''=ay''-bx'',\xi'''=ax'''+by''',\eta'''=ay'''-bx'''$, so findet man
\[
\begin{aligned}
\eta&=p-\xi'=p-\eta''=\xi'''\\
\xi-\eta&=\xi'+\eta'-p=\eta''-\xi''=p-\xi'''-\eta'''\\
p-\xi-\eta&=\xi'-\eta'=\xi''+\eta''-p=\eta'''-\xi'''
\end{aligned}
\]
woraus die Richtigkeit des Satzes mithilfe der Kriterien unmittelbar folgt. Da wiederum $r=m+ir'''$ gilt, sieht man leicht: Gehört $r$ nach Voraussetzung zu $C'$, so gehören die Zahlen $r',r'',r'''$ der Reihe nach zu $C'',C''',C$; gehört jene zu $C''$, so gehören diese zu $C''',C,C'$; und gehört jene schließlich zu $C'''$, so gehören diese zu $C,C',C''$.

Zugleich folgt daraus, dass die einzelnen Sammlungen $C,C',C'',C'''$ gleich viele Zahlen enthalten, nämlich $\frac14(p-1)$.

\par\medskip\Needspace{4\baselineskip}\pdfbookmark[1]{Art. 37}{paper-five-native151-art37}\centerline{\textbf{37.}}\nobreak\smallskip
\textsc{Satz.} \emph{Es bezeichne $k$ eine ganze Zahl, die durch $m$ nicht teilbar ist. Multipliziert man jede Zahl der Sammlung $C$ mit $k$ und verteilt die einfach kleinsten Reste der Produkte nach dem
\SourceLeaf{152}%
Modul $m$ auf die Sammlungen $C,C',C'',C'''$, wobei die Anzahlen der zu diesen einzelnen Sammlungen gehörigen Reste der Reihe nach mit $c,c',c'',c'''$ bezeichnet werden, so ist der Charakter der Zahl $k$ bezüglich des Moduls $m$ $\equiv c'+2c''+3c'''\pmod4$.}

\Needspace{22\baselineskip}
\emph{Beweis.} Es seien jene $c$ kleinsten Reste, die zu $C$ gehören, $\alpha,\beta,\gamma,\delta$ usw.; ferner seien die $c'$ Reste, die zu $C'$ gehören, $m+i\alpha',m+i\beta',m+i\gamma',m+i\delta'$ usw.; weiter seien die $c''$ Reste, die zu $C''$ gehören, $(1+i)m-\alpha'',(1+i)m-\beta'',(1+i)m-\gamma'',(1+i)m-\delta''$ usw.; und schließlich seien die $c'''$ Reste, die zu $C'''$ gehören, $im-i\alpha''',im-i\beta''',im-i\gamma''',im-i\delta'''$ usw. Wir betrachten nun vier Produkte, nämlich:
\begin{enumerate}
\renewcommand{\labelenumi}{\arabic{enumi})}
\item das Produkt aus allen $\frac14(p-1)$ Zahlen, die die Sammlung $C$ bilden:
\item das Produkt der Produkte, die durch Multiplikation jeder dieser Zahlen mit $k$ entstanden sind;
\item das Produkt aus den kleinsten Resten dieser Produkte, also aus den Zahlen $\alpha,\beta,\gamma,\delta$ usw., $m+i\alpha',m+i\beta'$ usw. usw.
\item das Produkt aus allen $c+c'+c''+c'''$ Zahlen $\alpha,\beta,\gamma,\delta$ usw., $\alpha',\beta',\gamma',\delta'$ usw., $\alpha'',\beta'',\gamma'',\delta''$ usw.
\end{enumerate}
Bezeichnen wir diese vier Produkte der Reihe nach mit $P,P',P'',P'''$, so gilt offenbar
\[
\begin{gathered}
P'=k^{\frac14(p-1)}P,\quad P'\equiv P'',\quad
P''\equiv P'''i^{c'+2c''+3c'''}\pmod m
\end{gathered}
\]
und daher
\[
Pk^{\frac14(p-1)}\equiv P'''i^{c'+2c''+3c'''}\pmod m
\]

Es ist aber leicht zu sehen, dass die Zahlen $\alpha',\allowbreak\beta',\allowbreak\gamma',\allowbreak\delta'$ usw., $\alpha'',\allowbreak\beta'',\allowbreak\gamma'',\allowbreak\delta''$ usw., $\alpha''',\allowbreak\beta''',\allowbreak\gamma''',\allowbreak\delta'''$ usw. sämtlich zur Sammlung $C$ gehören und sowohl untereinander als auch von den Zahlen $\alpha,\allowbreak\beta,\allowbreak\gamma,\allowbreak\delta$ usw. verschieden sind, ebenso wie diese selbst untereinander verschieden sind. Alle diese Zahlen zusammen müssen daher, wenn man von ihrer Reihenfolge absieht, mit allen Zahlen, die die Sammlung $C$ bilden, vollständig übereinstimmen. Daraus folgt $P=P'''$, also
\[
Pk^{\frac14(p-1)}\equiv Pi^{c'+2c''+3c'''}\pmod m
\]
Da schließlich keiner der einzelnen Faktoren des Produkts $P$ durch $m$ teilbar ist, folgt
\[
k^{\frac14(p-1)}\equiv i^{c'+2c''+3c'''}\pmod m
\]
sodass $c'+2c''+3c'''$ der Charakter der Zahl $k$ bezüglich des Moduls $m$ ist. Q.~E.~D.

\SourceLeaf{153}%
\Needspace{12\baselineskip}
\Article{72}
Um den allgemeinen Satz des vorigen Artikels auf die Zahl $1+i$ anzuwenden, müssen wir die Sammlung $C$ erneut in zwei kleinere Sammlungen $G$ und $G'$ unterteilen. Zur Sammlung $G$ rechnen wir die Zahlen $x+yi$, für die $ax+by=\xi$ kleiner als $\frac12p$ ist, zur anderen Sammlung $G'$ diejenigen, für die $\xi$ größer als $\frac12p$ ist. Die Anzahlen der in den Sammlungen $G,G'$ enthaltenen Zahlen bezeichnen wir der Reihe nach mit $g,g'$, sodass $g+g'=\frac14(p-1)$ gilt.

Das vollständige Kriterium für die zu $G$ gehörigen Zahlen besteht also darin, dass die drei Zahlen $\eta,\xi-\eta,p-2\xi$ positiv sind. Denn die dritte Bedingung für die Sammlung $C$, nach der $p-\xi-\eta$ positiv sein muss, ist in diesen Bedingungen bereits enthalten, da $p-\xi-\eta=(\xi-\eta)+(p-2\xi)$ gilt. Ebenso besteht das vollständige Kriterium für die zu $G'$ gehörigen Zahlen in positiven Werten der drei Zahlen $\eta,p-\xi-\eta,2\xi-p$.

Daraus folgt leicht, dass das Produkt einer beliebigen Zahl der Sammlung $G$ mit der Zahl $1+i$ zur Sammlung $C'''$ gehört. Setzt man nämlich
\[
\begin{gathered}
(x+yi)(1+i)=x'+y'i,\quad\text{und}\quad ax'+by'=\xi',\ ay'-bx'=\eta',\ \text{erhält man}\\
\xi'=\xi-\eta,\quad\eta'-\xi'=2\eta,\quad p-\xi'-\eta'=p-2\xi
\end{gathered}
\]
so ist das Kriterium dafür, dass die Zahl $x+yi$ zur Sammlung $G$ gehört, identisch mit dem Kriterium dafür, dass die Zahl $x'+y'i$ zur Sammlung $C'''$ gehört.

Auf genau dieselbe Weise zeigt man, dass das Produkt einer beliebigen Zahl der Sammlung $G'$ mit $1+i$ zur Sammlung $C''$ gehört.

\Needspace{12\baselineskip}
Geben wir also im vorigen Artikel $k$ den Wert $1+i$, so gilt $c=0,c'=0,c''=g',c'''=g$, und der Charakter der Zahl $1+i$ wird $3g+2g'=\frac12(p-1)+g$. Da die Charaktere der Zahlen $i,-1$ $\frac14(p-1),\frac12(p-1)$ sind, sind die Charaktere der Zahlen $-1+i,-1-i,1-i$ der Reihe nach $\frac34(p-1)+g,g,\frac14(p-1)+g$. Somit kommt es nun ganz auf die Bestimmung der Zahl $g$ an.
